% Part: sets-functions-relations
% Chapter: arithmetization
% Section: cuts
%
\documentclass[../../../include/open-logic-section]{subfiles}

\begin{document}

\olfileid{sfr}{arith}{cuts}
\olsection{Saka $\Rat$ menyang $\Real$}

Intine, Sipat Kajangkepan nuduhake yen saben titik $\alpha$ ing garis
real mbage garis iku dadi rong paro kanthi sampurna: paro kang duwe
$\alpha$ minangka wates ndhuwur paling cilik, lan paro kang duwe
$\alpha$ minangka wates ngisor paling gedhe. Kanggo \emph{mbangun}
wilangan real saka wilangan rasional, Dedekind ngusulake supaya kita
nganggep wilangan real minangka \emph{potongan} kang misahake wilangan
rasional. Tegese, kita ngenali $\sqrt{2}$ karo \emph{potongan} kang
misahake wilangan rasional $< \sqrt{2}$ saka wilangan rasional
$> \sqrt{2}$.

Ayo dirapekake. Yen wilangan rasional dipotong dadi rong paro, pamisahan
kang digawe bisa dikenali kanthi unik mung saka \emph{paro ngisor}.
Dadi, supaya luwih tliti, kita menehi definisi iki:

\begin{defn}[Potongan]
\emph{Potongan} $\alpha$ yaiku sembarang segmen wiwitan murni kang ora
kosong saka wilangan rasional lan ora duwe unsur paling gedhe. Tegese,
$\alpha$ iku potongan yen lan mung yen:
\begin{enumerate}
	\item \emph{ora kosong, murni}: $\emptyset \neq \alpha \subsetneq \Rat$
	\item \emph{wiwitan}: kanggo kabeh $p,q \in \Rat$: yen $p < q \in \alpha$ mula $p \in \alpha$
	\item \emph{tanpa maksimum}: kanggo saben $p \in \alpha$ ana $q \in \alpha$ saengga $p < q$
\end{enumerate}
Banjur $\Real$ iku himpunan kabeh potongan.
\end{defn}

Dadi saiki kita bisa kandha yen $\sqrt{2} = \Setabs{p \in \Rat}{p^2 < 2\text{
utawa }p < 0}$. Mesthi wae, kita kudu mriksa yen iki pancen
\emph{potongan}, nanging rinciane dipindhah menyang \olref[check]{sec}.

Kaya sadurunge, sawise netepake sawetara obyek, sabanjure kita kudu
netepake fungsi lan relasi dhasar ing obyek-obyek mau. Kita miwiti saka
kang gampang:
\begin{align*}
	\alpha \leq \beta \text{ yen lan mung yen }\alpha \subseteq \beta
\end{align*}
Definisi tatanan iki marakake kita bisa \emph{nyebutake} asil pokok,
yaiku himpunan potongan nduweni Sipat Kajangkepan. Yen ditulis kanthi
jangkep, pratelane mangkene. Yen $S$ iku himpunan potongan kang ora kosong
lan duwe wates ndhuwur, mula $S$ duwe wates ndhuwur paling cilik. Kanthi
luwih rinci: ana potongan $\lambda$ kang dadi wates ndhuwur kanggo $S$,
yaiku\ $(\forall \alpha \in S)\alpha \subseteq \lambda$, lan $\lambda$ iku
potongan paling cilik kang kaya mangkono, yaiku\ $(\forall \beta \in
\Real)((\forall \alpha \in S)\alpha \subseteq \beta \lif \lambda \subseteq
\beta)$. Saiki iki buktine:

\begin{thm}\ollabel{realcompleteness}
Himpunan potongan nduweni Sipat Kajangkepan.
\end{thm}

\begin{proof}
Upamakna $S$ iku sembarang himpunan potongan kang ora kosong lan duwe
wates ndhuwur. Netepna $\lambda = \bigcup S$.
%\Setabs{p \in \Rat}{(\exists \alpha \in S)p \in \alpha}$$
Kita luwih dhisik ngaku yen $\lambda$ iku potongan:
\begin{enumerate}
\item Amarga $S$ ora kosong, paling ora ana siji potongan ing $S$, mula
$\emptyset \neq \lambda$. Amarga $S$ iku himpunan potongan, $\lambda
\subseteq \Rat$. Amarga $S$ duwe wates ndhuwur, ana $p \in \Rat$ kang
ora ana ing saben potongan $\alpha \in S$. Dadi $p\notin \lambda$, mula
$\lambda \subsetneq \Rat$. Cathetan koreksi sumber OLP-0045: anane
potongan ing himpunan kasebut dumadi amarga himpunane ora kosong, dudu
amarga duwe wates ndhuwur kaya kang kasebut ing sumber Inggris.
\item Upamakna $p < q \in \lambda$. Dadi ana $\alpha \in S$ saengga
$q \in \alpha$. Amarga $\alpha$ iku potongan, $p \in \alpha$. Dadi
$p \in \lambda$.
\item Upamakna $p \in \lambda$. Dadi ana $\alpha \in S$ saengga
$p \in \alpha$. Amarga $\alpha$ iku potongan, ana $q \in \alpha$
saengga $p < q$. Dadi $q \in \lambda$.
\end{enumerate}
Iki mbuktekake pratelan mau. Kajaba iku, cetha yen $(\forall \alpha \in
S)\alpha \subseteq \bigcup S = \lambda$, yaiku\ $\lambda$ dadi wates
ndhuwur kanggo $S$. Saiki upamakna $\beta \in \mathbb{R}$ uga dadi wates
ndhuwur, yaiku\ $(\forall \alpha \in S)\alpha \subseteq \beta$. Kanggo
sembarang $p \in \Rat$, yen $p \in \lambda$, mula ana $\alpha \in S$
saengga $p \in \alpha$, mula $p \in \beta$. Kanthi nggeneralisasi,
$\lambda \subseteq \beta$. Dadi $\lambda$ iku wates ndhuwur
\emph{paling cilik} kanggo $S$.
\end{proof}

Dadi kita duwe akeh obyek kang nyukupi Sipat Kajangkepan. Salah siji
cara kanggo ngandharake iki yaiku: ora ana ``longkangan'' ing
potongan-potongan kita. (Dadi, nggawe ``potongan'' sabanjure saka wilangan
real, dudu saka wilangan rasional, ora bakal ngasilake obyek anyar kang
narik kawigaten.)

Sabanjure, kita kudu netepake sawetara operasi ing wilangan real. Kita
miwiti kanthi nyisipake wilangan rasional menyang wilangan real kanthi
netepake $p_\Real = \Setabs{q \in \Rat}{q < p}$ kanggo saben $p \in \Rat$.
Banjur kita netepake:
\begin{align*}
	\alpha + \beta &= \Setabs{p + q}{p \in \alpha \land q \in \beta}\\
%	\alpha - \beta &\defis \Setabs{p - q}{p \in \alpha \land q \in \Rat \setminus \beta}\\
	\alpha \times \beta &=
	\Setabs{p \times q}{0 \leq p \in \alpha \land 0 \leq q \in \beta} \cup 0^\mathbb{R} & \text{yen }\alpha, \beta \geq 0_\Real
%	\alpha \div \beta &\defis
%\Setabs{\nicefrac{p}{q}}{p \in \alpha \land q \in \Rat \setminus \beta}  & \text{yen }\alpha \geq 0_\Real, \beta > 0_\Real
\end{align*}
Kanggo nangani kasus ping-pingan liyane, luwih dhisik netepna: %yen $0_\Real \times \alpha = 0_\Real = \alpha \times 0_\Real$, lan tambahna:
\begin{align*}
	-\alpha &= \Setabs{p - q}{p < 0 \land q \notin \alpha}
\end{align*}
banjur netepna:
\begin{align*}
	\alpha \times \beta &\defis
	\begin{cases}
		\mathord{-}\alpha \times \mathord{-}\beta &\text{yen }\alpha < 0_\Real\text{ lan }\beta < 0_\Real\\
		\mathord{-}(\mathord{-}\alpha \times \beta) &\text{yen }\alpha < 0_\Real \text{ lan }\beta > 0_\Real\\
		\mathord{-}(\alpha \times \mathord{-}\beta) &\text{yen }\alpha > 0_\Real \text{ lan }\beta < 0_\Real
	\end{cases}
\end{align*}
Banjur kita kudu mriksa yen saben definisi iki tansah ngasilake potongan.
Pungkasan, kita kudu nindakake pamriksa kang gampang (nanging dawa) yen
potongan-potongan kang ditetepake iki tumindak persis kaya kang kudune.
Nanging kabeh rinciane dipindhah menyang \olref[check]{sec}.
\end{document}
